% kyber (sk07) -- infrastructure hand-off report
%
% BUILD: this document embeds the topology diagram with \includesvg, which needs
%        Inkscape on PATH and shell-escape enabled:
%            cd report && latexmk -pdf -shell-escape kyber-report.tex
%        If Inkscape is unavailable, pre-convert the SVG to PDF once
%        (e.g. `inkscape ../kyber-network-topology.svg --export-filename=topology.pdf`
%        or `rsvg-convert -f pdf -o topology.pdf ../kyber-network-topology.svg`)
%        and switch the figure to the commented \includegraphics fallback in section 2.
%
\documentclass[11pt,a4paper]{article}

\usepackage[T1]{fontenc}
\usepackage[utf8]{inputenc}
\usepackage[margin=1in]{geometry}
\usepackage{booktabs}
\usepackage{tabularx}
\usepackage{array}
\usepackage{xcolor}
\usepackage{enumitem}
\usepackage{svg}            % \includesvg -- needs Inkscape + --shell-escape (see header)
\usepackage[hidelinks]{hyperref}

\newcommand{\code}[1]{\texttt{#1}}
\setlist{nosep,leftmargin=1.4em}
\renewcommand{\arraystretch}{1.15}

\title{\textbf{Kyber (SK07)} \\[2pt] \large Infrastructure Overview \& Hand-off Report}
\author{Group SK07}
\date{June 2026}

\begin{document}
\maketitle


\tableofcontents
\bigskip
\hrule

\section{Introduction}

kyber is the in-house network for a fictional start-up: a complete company network and
service stack with segmented users and servers, dual-stack IPv4 and IPv6 throughout, an
on-prem user directory, a TLS-protected REST service backed by a database, monitoring,
remote VPN access, a RAFT-based highly-available service, and a documented firewall.
Everything is hosted on the team's own ESXi infrastructure --- there are no cloud-hosted
services.

Two ideas are worth absorbing before anything else, because they shape day-to-day work:

\begin{itemize}
  \item \textbf{Management is VPN-only.} The router does not accept SSH from the
        internet. To administer any box you first bring up the VPN
        (Section~\ref{sec:vpn}); the router then doubles as an SSH \emph{jump host} to
        reach DMZ and internal machines that have no public address
        (\code{ssh -J vyos@<router> <user>@<host>}).
  \item \textbf{Names resolve differently inside and outside (split DNS).} Inside the
        network, \code{api.kyber.local} and friends resolve to \emph{private} addresses
        and traffic never leaves the LAN; from outside, the same name reaches the one
        public address. This is intentional and is the reason most testing is done from
        the internal workstations.
\end{itemize}

The internal domain is \code{kyber.local}. The ISP (the lab, ``LRK'') assigned one public
IPv4 address and one IPv6 \code{/62} block.

\section{Architecture at a Glance}
\label{sec:arch}

The network is built around a single VyOS 1.4.4 router, \code{kyber-rtr}, which owns four
network interfaces --- one per segment --- and enforces all routing, NAT, and firewalling
between them. Figure~\ref{fig:topology} shows the segments and how they connect to the
outside world.

\begin{figure}[htbp]
  \centering
  % Requires Inkscape on PATH and compilation with --shell-escape (see file header).
  \includesvg[width=\linewidth, inkscapelatex=false]{../kyber-network-topology.svg}
  % Fallback without Inkscape: pre-convert the SVG to topology.pdf and use instead:
  % \includegraphics[width=\linewidth]{topology.pdf}
  \caption{kyber sk07 network topology: the VyOS router \code{kyber-rtr} joins the WAN,
           the internal user segment, the DMZ server segment, and the IPv6-only segment.}
  \label{fig:topology}
\end{figure}

The design follows a classic segmentation: \textbf{users and servers are kept apart}. End
users live on the \emph{internal} segment; all servers live in the \emph{DMZ}; a dedicated
\emph{IPv6-only} segment demonstrates ULA addressing with prefix translation. The router
treats each interface as a firewall \emph{zone} (\code{WAN}, \code{INTERNAL}, \code{DMZ},
\code{V6ONLY}), plus \code{LOCAL} for the router itself and \code{VPN} for remote clients.

\begin{table}[htbp]
\centering
\footnotesize
\setlength{\tabcolsep}{4pt}
\begin{tabularx}{\linewidth}{@{}l l l l l X@{}}
\toprule
\textbf{VM} & \textbf{OS} & \textbf{Segment} & \textbf{IPv4} & \textbf{IPv6} & \textbf{Roles} \\
\midrule
\code{kyber-rtr}    & VyOS 1.4.4    & all four   & per-iface (Tbl.~\ref{tab:addr}) & per-iface & Router, firewall, NAT/NPTv6, DHCP, DNS forwarder, NTP, VPN, SNMP, NetFlow export \\
\code{kyber-app-01} & Ubuntu Server & DMZ        & 192.168.7.10  & \code{::10}  & REST API, PostgreSQL (Patroni primary), etcd, nginx, keepalived MASTER \\
\code{kyber-app-02} & Ubuntu Server & DMZ        & 192.168.7.11  & \code{::11}  & REST API, PostgreSQL (Patroni standby), etcd, nginx, keepalived BACKUP \\
\textit{API VIP}    & ---           & DMZ        & 192.168.7.100 & \code{::100} & Floating service address; \code{api.kyber.local} resolves here \\
\code{kyber-mon}    & Ubuntu Server & DMZ        & 192.168.7.20  & \code{::20}  & Prometheus, Grafana, ntopng, etcd witness \\
\code{kyber-ldap}   & AlmaLinux 10  & DMZ        & 192.168.7.30  & \code{::30}  & FreeIPA: LDAP + Kerberos + CA + authoritative DNS \\
\code{kyber-ws-01}  & Ubuntu Desktop& internal   & dynamic       & dynamic      & Linux end-user workstation \\
\code{kyber-ws-02}  & Windows Server& internal   & dynamic       & dynamic      & Windows end-user workstation \\
\code{kyber-ipv6}   & Ubuntu Server & IPv6-only  & --- none ---  & SLAAC ULA    & Demonstrates the IPv6-only + NPTv6 segment \\
\bottomrule
\end{tabularx}
\caption{Virtual-machine inventory. The IPv6 suffix column is under the DMZ prefix
         \code{2001:1470:fffd:99::/64}; internal workstations lease from
         \code{2001:1470:fffd:9a::/64}. Heterogeneous operating systems (VyOS, Ubuntu,
         AlmaLinux, Windows) are used by design.}
\label{tab:vms}
\end{table}

All DMZ servers receive \textbf{fixed} addresses, but via DHCP \emph{reservations}
(MAC-to-IP bindings), not hand-configured static addresses --- so a rebuilt VM always comes
back on the same address without per-host editing. Internal workstations take ordinary
dynamic leases.

\section{Addressing \& Core Router Services}
\label{sec:addr}

The assigned IPv6 block \code{2001:1470:fffd:98::/62} splits into exactly four \code{/64}
networks: the first for the internet uplink, two for the
dual-stack LAN segments, and the last as the outer prefix for the IPv6-only segment's
prefix translation. Table~\ref{tab:addr} is the master addressing plan.

\begin{table}[htbp]
\centering
\footnotesize
\setlength{\tabcolsep}{5pt}
\begin{tabularx}{\linewidth}{@{}l l l l X@{}}
\toprule
\textbf{Segment} & \textbf{IPv4 (router IP)} & \textbf{IPv6 \code{/64} (router IP)} & \textbf{Iface} & \textbf{IPv6 auto-config} \\
\midrule
WAN / uplink & 88.200.24.237/25 (GW .129) & \code{2001:1470:fffd:98::2} (GW \code{::1}) & eth0 & static \\
Internal     & 10.7.0.0/24 (\code{.1})    & \code{2001:1470:fffd:9a::/64} (\code{::1})  & eth1 & DHCPv6 (stateful) \\
DMZ          & 192.168.7.0/24 (\code{.1}) & \code{2001:1470:fffd:99::/64} (\code{::1})  & eth2 & DHCPv6 (stateful) \\
IPv6-only    & --- none ---               & ULA \code{fd07:1:1:1::/64} (\code{::1})     & eth3 & SLAAC \\
\bottomrule
\end{tabularx}
\caption{Master address plan. The NPTv6 \emph{outer} prefix \code{2001:1470:fffd:9b::/64}
         is not assigned to any interface; it is the public face of the IPv6-only ULA.}
\label{tab:addr}
\end{table}

The router runs the usual ``general settings'' a company needs, all centralised on
\code{kyber-rtr}:

\begin{description}[leftmargin=1.4em,style=nextline]
  \item[NAT44 (outbound internet)] Both LAN subnets are source-masqueraded behind the
    single public IPv4 on egress, so users and servers reach the internet while presenting
    one address.
  \item[NPTv6 (IPv6 prefix translation)] The IPv6-only segment uses private ULA addresses
    internally; on the way out, the router rewrites the ULA prefix
    \code{fd07:1:1:1::/64} to the public \code{2001:1470:fffd:9b::/64} (RFC~6296), and the
    reverse on the way back. This gives IPv6-only hosts internet reach without NAT44-style
    port translation.
  \item[DHCP / DHCPv6 / SLAAC] IPv4 is served by DHCP on both LANs (a dynamic pool
    \code{10.7.0.100--200} internally; reservations in the DMZ). For IPv6, both dual-stack
    LANs use \textbf{stateful DHCPv6} (router advertisements set the \emph{managed} flag and
    suppress autonomous addressing), while \textbf{SLAAC} is used on the IPv6-only segment
    (eth3 advertises its ULA prefix autonomously).
  \item[DNS forwarding \& split DNS] The router is the resolver for every internal segment.
    Queries for \code{kyber.local} (and the DMZ reverse zone) are forwarded to the internal
    FreeIPA DNS at \code{192.168.7.30}, which returns private addresses; everything else
    goes to public upstream resolvers.
  \item[NTP relay] The router synchronises to public pools and re-serves time to the
    internal, DMZ, IPv6-only, and VPN clients.
  \item[SNMP agent] A read-only SNMPv2c agent (community \code{kyber-ro}) exposes interface
    counters and system metrics, \textbf{restricted to the monitoring host only}
    (\code{kyber-mon}) on both stacks.
  \item[NetFlow export] All four interfaces export NetFlow~v9 flow records to
    \code{kyber-mon} for traffic analysis (Section~\ref{sec:mon}).
\end{description}

\section{Firewall \& Security Posture}
\label{sec:fw}

The firewall is \textbf{zone-based and fully dual-stack}: every rule is written for both
IPv4 and IPv6, since the network is dual-stack end to end. The default action for every
zone pair is \textbf{drop}; a single global stateful rule accepts established/related return
traffic and drops invalid packets, so each ruleset only has to list the \emph{new}
connections it intends to permit. Table~\ref{tab:fw} is the full zone-pair matrix; the
intent is summarised first.

\begin{itemize}
  \item \textbf{From the internet (WAN), almost nothing is allowed in.} Only the \code{WAN}
        rows of Table~\ref{tab:fw} are reachable. The internal segment is never reachable
        from the WAN; the DMZ is reachable only on \code{tcp/443}.
  \item \textbf{Users (internal) reach out freely} to the internet, and reach the DMZ only
        on the specific service ports they need (HTTPS, SSH for admins, DNS, LDAPS).
  \item \textbf{Servers (DMZ) are constrained.} They may reach the internet for updates and
        time, but they \textbf{cannot initiate connections back to the user segment} --- if
        a server were compromised, this blocks lateral movement to user machines. That drop
        is logged.
  \item \textbf{IPv6 hygiene.} WAN ingress drops spoofed/bogon sources, and crucially drops
        inbound IPv6 Router Advertisements, so an outside party cannot hijack the default
        route of LAN hosts (rogue-RA protection). The neighbor-discovery and error ICMPv6
        types that IPv6 genuinely needs are permitted.
\end{itemize}

\begin{table}[htbp]
\centering
\footnotesize
\setlength{\tabcolsep}{4pt}
\begin{tabularx}{\linewidth}{@{}l >{\hsize=.95\hsize}X >{\hsize=1.05\hsize}X@{}}
\toprule
\textbf{From $\rightarrow$ To} & \textbf{Allowed (IPv4 + IPv6)} & \textbf{Purpose} \\
\midrule
WAN $\rightarrow$ LOCAL      & ICMP echo (rate-limited); \code{udp/1194} & Ping; OpenVPN. No SSH from the internet. \\
WAN $\rightarrow$ DMZ        & \code{tcp/443}                            & Published HTTPS: REST API + Grafana/ntopng dashboards. \\
WAN $\rightarrow$ INTERNAL   & --- (drop, logged)                        & User segment never reachable from the internet. \\
WAN $\rightarrow$ V6ONLY     & --- (drop)                                & No inbound to IPv6-only hosts. \\
INTERNAL $\rightarrow$ WAN   & everything                                & Users' unrestricted internet access. \\
INTERNAL $\rightarrow$ DMZ   & \code{tcp/443,22,636}; \code{tcp+udp/53}; \code{tcp/80}$\to$CA & Services, SSH admin, DNS, LDAPS, revocation checks. \\
INTERNAL $\rightarrow$ LOCAL & \code{tcp/22}; \code{tcp+udp/53}; \code{udp/123,67}; ICMP echo & Router mgmt, DNS, NTP, DHCP, ping. \\
DMZ $\rightarrow$ WAN        & \code{tcp/80,443}; \code{udp/123}; ICMP echo & OS/package updates, NTP, diagnostics. \\
DMZ $\rightarrow$ INTERNAL   & --- (drop, logged)                        & Servers cannot initiate to users (limits lateral movement). \\
DMZ $\rightarrow$ LOCAL      & \code{tcp+udp/53}; \code{udp/123}; \code{udp/161} (mon only); ICMP echo & DNS, NTP, SNMP polling. \\
V6ONLY $\rightarrow$ WAN     & everything (IPv6)                         & IPv6-only internet egress via NPTv6. \\
V6ONLY $\rightarrow$ LOCAL   & \code{tcp+udp/53}; \code{udp/123}; ICMPv6 & DNS, NTP, neighbour discovery. \\
VPN $\rightarrow$ LOCAL      & \code{tcp/22}; \code{tcp+udp/53}; \code{udp/123} & Router mgmt, DNS, NTP. \\
VPN $\rightarrow$ INTERNAL   & \code{tcp/22,3389}; ICMP echo             & SSH to Linux, RDP to Windows, ping. \\
VPN $\rightarrow$ DMZ        & \code{tcp/443,22,636}                     & Services, SSH, LDAPS. \\
LOCAL $\rightarrow$ any      & everything                                & Router-initiated traffic (NTP, DNS, updates). \\
\bottomrule
\end{tabularx}
\caption{Zone-pair firewall matrix; the same policy is encoded for IPv4 and IPv6. Any pair
         not listed defaults to \textbf{drop}; security-relevant drops (WAN inbound,
         DMZ $\rightarrow$ INTERNAL) are logged. Return traffic for established connections is
         always accepted by a global stateful rule.}
\label{tab:fw}
\end{table}

\section{Remote-Access VPN}
\label{sec:vpn}

Remote access is provided by \textbf{OpenVPN}, run natively on the VyOS router (no separate
appliance). It listens on \code{udp/1194} and gives connected clients a dual-stack tunnel
(\code{10.7.99.0/24} and \code{fd07:99::/64}). It is a \textbf{split tunnel}: only routes to
the company's own LAN and DMZ subnets are pushed, so a remote worker's general internet
traffic does not transit the company.

Authentication is tied to the directory: users log in with their \textbf{FreeIPA username
and password}, validated live over LDAPS, and are authorised by membership of the
\code{vpn-users} group. There are no per-user keys or certificates to issue or revoke ---
\textbf{off-boarding is a single directory action} (disable the user or remove them from
the group, and their next connection fails). The tunnel's control channel is secured by its
own small certificate authority generated on the router; note this is a \emph{separate} CA
from the FreeIPA one used everywhere else (Section~\ref{sec:idm}). Because router SSH is
closed to the WAN, the VPN is the front door for all administration.

\section{Identity \& DNS --- FreeIPA}
\label{sec:idm}

The user directory is \textbf{FreeIPA} on \code{kyber-ldap}, chosen because it bundles, in
one product, everything the rest of the build leans on: an \textbf{LDAP} directory,
\textbf{Kerberos}, an internal \textbf{certificate authority}, and an integrated
\textbf{authoritative DNS} (BIND) for \code{kyber.local}.

\begin{itemize}
  \item \textbf{DNS.} FreeIPA is authoritative for \code{kyber.local} and the DMZ reverse
        zone, holding the A/AAAA records for every server and the service names
        (\code{api}, \code{grafana}, \code{ntopng}). The router forwards internal queries
        here --- this is the ``private'' side of split DNS.
  \item \textbf{Certificate authority.} The FreeIPA CA issues the TLS certificates for
        \emph{all} internal services (the REST API, Grafana, ntopng, LDAPS, and the etcd /
        PostgreSQL replication layer). Certificates are requested and auto-renewed by
        \code{certmonger} on each host, so nothing is hand-managed or self-signed. Every
        client trusts this one CA, which is the single setup step needed for internal HTTPS
        to validate.
\end{itemize}

\begin{table}[htbp]
\centering
\small
\begin{tabular}{@{}l l l@{}}
\toprule
\textbf{User} & \textbf{Groups} & \textbf{Purpose} \\
\midrule
\code{alice} & \code{admins}, \code{vpn-users}      & Administrator; can VPN in \\
\code{bob}   & \code{vpn-users}, \code{ipausers}    & Ordinary VPN user \\
\code{carol} & \code{api-writers}                   & Authorised to call REST write endpoints \\
\code{dave}  & \code{ipausers}                      & Plain user --- the ``denied'' test case \\
\code{luka}, \code{urban} & \code{vpn-users}        & Real team members' VPN accounts \\
\bottomrule
\end{tabular}
\caption{Directory users and the groups that gate access. \code{vpn-users} controls VPN
         login; \code{api-writers} controls REST write access. \code{dave} belongs to
         neither, which makes him the standing negative test for both gates.}
\label{tab:users}
\end{table}

\section{The REST Service}
\label{sec:api}

The customer-facing service is a \textbf{REST API} built with FastAPI (Python) over
SQLAlchemy and \textbf{PostgreSQL}. It exposes two related resources in a classic
one-to-many relationship --- \textbf{customers} and their \textbf{orders} (deleting a
customer cascades to their orders) --- with full create/read/update/delete operations on
each.

Three behaviours are worth highlighting:

\begin{description}[leftmargin=1.4em,style=nextline]
  \item[Content negotiation] The same resource is returned as \textbf{JSON}, \textbf{XML},
    or an \textbf{HTML} table, selected by the request's \code{Accept} header --- three
    formats from one code path.
  \item[Authentication \& authorisation] Read operations (\code{GET}) are public. Write
    operations (\code{POST}/\code{PUT}/\code{DELETE}) require HTTP Basic credentials that
    \textbf{bind successfully to FreeIPA over LDAPS} \emph{and} belong to the
    \code{api-writers} group. So \code{carol} can write, \code{dave} is forbidden (403),
    and an unauthenticated write is challenged (401).
  \item[Transport] \textbf{nginx} terminates TLS in front of the application on
    \code{443} with a FreeIPA-issued certificate, offers \textbf{HTTP/2} alongside
    HTTP/1.1, and listens on \textbf{both IPv4 and IPv6}.
\end{description}

The database password and directory settings are supplied from an environment file on the
host, never embedded in the application code.

\section{High Availability \& the Data Tier (RAFT)}
\label{sec:ha}

The service is designed to survive the loss of a node, at both the web tier and the data
tier.

\textbf{Web tier.} Two application nodes (\code{kyber-app-01}, \code{kyber-app-02}) each run
the API behind nginx. A \textbf{keepalived} virtual IP (\code{192.168.7.100} /
\code{2001:1470:fffd:99::100}) floats between them, and \code{api.kyber.local} resolves to
that VIP. The VIP-holder's nginx load-balances \textbf{active-active} across both backends,
so a single node or process failure does not interrupt service --- the VIP moves (or the
dead backend is simply skipped) within seconds.

\textbf{Data tier (RAFT).} A three-node \textbf{etcd} cluster runs on \code{kyber-app-01},
\code{kyber-app-02}, and \code{kyber-mon} (the third acts purely as a voting
\emph{witness}, so losing one application node never costs quorum). etcd is the RAFT
consensus store; its consumer is \textbf{Patroni}, which manages PostgreSQL as an
\textbf{auto-failover primary plus hot standby} with streaming replication. Patroni records
leadership in etcd, so there is never a split brain, and if the primary node fails the
standby is promoted automatically. The API connects with a multi-host database URL and
always lands on whichever node is currently primary --- no manual intervention. All of this traffic is
internal to the DMZ and secured with FreeIPA-issued certificates.

\section{Publishing Services to the Outside World}
\label{sec:publish}

The company owns exactly one public IPv4 address but serves several named HTTPS services
(\code{api}, \code{grafana}, \code{ntopng}), so the edge works differently per stack:

\begin{itemize}
  \item \textbf{IPv4} --- the router performs \textbf{destination NAT}: inbound
        \code{88.200.24.237:443} is forwarded to the internal VIP. On the VIP, a lightweight
        \textbf{SNI edge} (nginx in TLS-passthrough mode) reads the requested hostname from
        the TLS handshake and forwards the still-encrypted connection to the matching
        backend, so one address and one port serve all three names with end-to-end TLS
        preserved (no certificates are moved to the edge).
  \item \textbf{IPv6} --- no NAT is needed: every server already has a globally routable
        address, so each service's name points straight at it and the router simply routes.
\end{itemize}

There is no registered public domain, so external clients map the three service names to the
public address with a local hosts file.

\section{Monitoring \& Flow Analysis}
\label{sec:mon}

Monitoring lives on \code{kyber-mon} and has two complementary halves:

\begin{itemize}
  \item \textbf{Metrics --- Prometheus + Grafana.} Prometheus scrapes at a short interval:
        a \code{node\_exporter} on every Linux host (CPU, memory, disk) and an
        \code{snmp\_exporter} that polls the router's interface counters via the
        \code{kyber-ro} community. \textbf{Grafana} presents this --- per-interface WAN
        throughput over time, and host CPU/memory --- published over HTTPS at
        \code{grafana.kyber.local} with a FreeIPA certificate.
  \item \textbf{Flows --- ntopng.} The router's NetFlow~v9 export is received by a collector
        and rendered by \textbf{ntopng} (top-talkers, per-flow byte counts), published the
        same way at \code{ntopng.kyber.local}. This gives the network administrator a view
        of \emph{who is talking to whom}, complementing Prometheus's \emph{how loaded is
        each box}.
\end{itemize}

Both dashboards are served over HTTPS with FreeIPA certificates and are reachable by name,
like the API (Section~\ref{sec:publish}).

\section{Clients \& the IPv6-only Segment}
\label{sec:clients}

Two end-user workstations on the internal segment run deliberately \textbf{different
operating systems}: \code{kyber-ws-01} (Ubuntu Desktop, Linux) and \code{kyber-ws-02}
(Windows Server with Desktop Experience). Both take dynamic leases, trust the FreeIPA CA, and
are used to verify the REST API end to end (content negotiation, HTTP/2, IPv6 reachability,
and the write-authorisation cases).

The \textbf{IPv6-only segment} carries \code{kyber-ipv6}, a host with \emph{no IPv4 stack at
all}. It auto-configures a ULA address by SLAAC and reaches the internet purely over IPv6
through NPTv6 (Section~\ref{sec:addr}).

\end{document}
